\clearpage
\fancyhead[L]{\small THE SUN’S GREATEST ALTITUDE.}
\hypertarget{native-465-sunheight}{}
\begin{center}{\large THE SUN’S GREATEST ALTITUDE.}\end{center}

\begin{center}
\rule{0.55\linewidth}{0.4pt}\\[2pt]
Astronomische Hülfstafeln for 1827, edited by \textsc{Schumacher}, p.~92.\\[2pt]
\rule{0.55\linewidth}{0.4pt}
\end{center}

As is well known, the Sun’s changing declination means that it reaches its greatest altitude away from the meridian. In the ascending signs this occurs \emph{after} its passage through the meridian and in the descending signs \emph{before} it. The instant of greatest altitude is noteworthy because, as Councillor \textsc{Gauss} ingeniously observed, it saves the calculator the trouble of allowing for the change in declination over the observation period when working with altitudes of
\SourceLeaf{466}the Sun measured near the meridian. The hour angles are the arguments of the corrections used to reduce altitudes measured away from noon to the noon altitude. The calculator need only reckon these angles from the instant when the Sun reaches its greatest altitude instead of from true noon as usual. This simple change in the zero point of the hour angles accounts for the effect of the change in declination over the entire series of observations. --- The following table gives, for every tenth day and for all altitudes of the pole from Gibraltar to Petersburg (from $36^\circ$ to $60^\circ$), the amount to add to or subtract from the instant of true noon to obtain the instant when the Sun reaches its greatest altitude. \textsc{Schumacher.}
